\documentclass[11pt,a4paper]{article}
\usepackage[margin=25mm]{geometry}
\usepackage{amsmath,amssymb,amsthm,booktabs,microtype}
\usepackage[T1]{fontenc}
\usepackage{lmodern}
\usepackage[colorlinks=true,linkcolor=blue,urlcolor=blue,citecolor=blue]{hyperref}
\newcommand{\FS}{\operatorname{FS}}
\newcommand{\Z}{\mathbb Z}
\newcommand{\T}{\mathcal T}
\newtheorem{theorem}{Theorem}[section]
\newtheorem{lemma}[theorem]{Lemma}
\newtheorem{proposition}[theorem]{Proposition}
\title{Subset Sums in a Fourfold Sumset\\
\large A checked continuation from the upper bound 115 to 103}
\author{Sum-of-Subset research project}
\date{14 October 2026}
\begin{document}
\maketitle
\begin{abstract}
Let $\FS(A)$ include the empty subset sum, and let $4A$ allow repeated
summands. Starting with the previously audited bound $N\le115$ for
$\FS(A)\subseteq4A$, we exclude every cardinality from 104 through 115
by one common argument. The resulting independently checked continuation
gives $71\le N\le103$. Its finite part consists of 24 explicit hole
witnesses and exact scalar inequalities, verified by a short standard-library
Python program. It uses no numerical optimization or search over candidate
sets. The structural lemmas are proved here, with the audited 115 foundation
explicitly identified. The repository's earlier claim $N\le84$ is stronger;
this report does not claim a new project record or certify that entire chain.
\end{abstract}
\tableofcontents

\section{Scope, prerequisites, and the main conclusion}
For a finite real set $X$, put
\[
 \FS(X)=\left\{\sum_{x\in Y}x:Y\subseteq X\right\},\qquad
 R(X)=\{x+y:x,y\in X,\ x\ne y\},\quad
 \Delta(X)=|2X|-|R(X)|.
\]
Here $2X=X+X$, $T_r=r(r+1)/2$, and $N$ is the largest feasible
cardinality. The following theorem concerns all real sets, not only integer
sets or sets of a prescribed shape.
\begin{theorem}\label{thm:main}
If $\FS(A)\subseteq4A$, then $|A|\le103$.
\end{theorem}
The 71-element construction from the accompanying research report remains
unchanged:
\[
 A=[-31,31]_{\Z}\cup\{\pm93,\pm155,\pm217,\pm713\}.
\]
The precise maximum is not determined here.

\subsection{The audited starting point}
This is a continuation of the 115 foundation, not a replacement for its
proof. We use the following previously checked facts, as stated in Part I
of \texttt{upper86\_complete\_20261005.tex}, with the peeling correction in
\texttt{proof\_review\_20261014.md}.
\begin{enumerate}
\item A feasible set remains feasible on adjoining zero. Any feasible real
set has an integer realization of the same cardinality and sign order,
preserving zero when present. To see the latter, fix one four-term
representation for every subset sum. These are homogeneous rational linear
equations; rational points are dense in their real kernel. Preserve the
strict order inequalities and then clear denominators.
\item Every feasible set has at most 115 elements.
\item Remove the four largest positive and four most negative elements,
taking all of a sign if fewer than four exist. With remainder $D$, $d=|D|$,
\begin{equation}\label{eq:cover}
 |\FS(D)|\le2|2D|+18d+85.
\end{equation}
The refined count consists of two leading $2D$ translates, eighteen
ordinary units of capacity $d$, at most 84 nonzero fixed contributions,
and zero. Its 110 fixed labels and 26 disjoint nonpositive groups are
also reproduced and checked in the newer research report and
\texttt{verification/verify.py}.
\item A four-extreme peel, preserving zero, gives
\begin{equation}\label{eq:peel}
 |\FS(F)|\ge|\FS(F')|+|R(F)|-1.
\end{equation}
When both signs suffice, remove two elements at each extreme. If the
extreme deleted sums are $\alpha,\beta$, the residue-class argument counts
the classes of $\alpha+\FS(F')$ modulo $\beta-\alpha$, not the original
classes of $\FS(F')$. In each translated class the two extreme translates
provide one additional point. Restricted sums in other classes then prove
\eqref{eq:peel}. This is the correction used throughout this report.
\item For $X_i=\{0<x_1<\cdots<x_i\}$, write
\[
 \mu(X_i)=|\FS(X_i)|-1-T_i=\sum_{j=1}^i y_j,\qquad y_j\ge0,
\]
\[
 \delta_0(X_i)=|\{x\in X_i:2x\notin R(X_i\cup\{0\})\}|,
 \quad\rho(X_i)=\mu(X_i)+2-2\delta_0(X_i).
\]
Then $\Phi=2-\rho=2\delta_0-\mu\le10$, and
$\rho(X_i)\ge\rho(X_{i-1})+y_i-2$. If $i\ge6$ and $y_i\le i-5$,
the small-increment inverse lemma gives
\begin{equation}\label{eq:small-inverse}
 X_i=t([1,h]_{\Z}\setminus D_0),\quad i\le h\le i+y_i,
 \quad |D_0|=h-i,\quad h\notin D_0.
\end{equation}
The circle-group argument proving this lemma uses the subset-sum growth
results of Balandraud--Girard--Griffiths--Hamidoune \cite{BGGH}.
Equality $y_i=0$ yields an initial progression; for $i\ge6$, $y_i=0,1$
implies $\Phi(X_i)\le2$.
\end{enumerate}
We additionally use the following precise external results. The stability
theorem of Carpenter--Defant--Kravitz \cite[Theorem 1.1]{CDK} says that for
$r\ge4$ and $0\le M\le r-4$, $\mu(X)\le M$ implies
$X=tB$, where $B$ is a positive integer set of size $r$ and
$\sum B\le T_r+M$. In particular a non-progression subset-sum support
has $\mu\ge r-3$. The two-set inverse theorem and central-interval theorem
are stated at their actual applications below. We do not assume any
post-115 endpoint exclusion from the archived reports.

\section{Positive-side defect estimates}
\begin{lemma}[Deleted intervals]\label{lem:deleted}
If $X=[1,r+q]_{\Z}\setminus D$, $|D|=q$, $r+q\notin D$, and
$r\ge\max(5,4q+1)$, then $\Phi(X)\le2$. For $q\ge1$ one also has
$\delta_0(X)\le2q$. The latter defect bound remains valid for $q=1,2$
and $r\ge7$.
\end{lemma}
\begin{proof}
The case $q=0$ is immediate. Write $h=r+q$ and
$C=\{h-d:d\in D\}$. The integer mass excess is
\[
 e_0=\sum X-T_r=\sum_{c\in C}c-T_{q-1}.
\]
There are $\min(k,h-k)$ pairwise disjoint representations of $2k$ by
distinct members of $[0,h]_{\Z}$. Thus a defective present $k$ must be
in $[1,q]$ or $[r,h]$. At most $q+1$ high defects occur. For $q\ge2$,
all $1,\ldots,q$ cannot be both present and defective: defectivity of 1
forces 2 absent. Hence at most $q-1$ low defects occur. For $q=1$,
the only extra defect besides the maximum can arise from deleting 2 or
$h-2$; these are distinct for $r\ge5$. This proves the defect counts.

Since $r>q$, $X$ is primitive. If $\mu\le r-4$, stability implies
$\mu=e_0$: every integer normalization of a primitive integer set is an
integer dilation, while $|\FS(X)|\le\sum X+1$. With no low defects, let
the high defects other than $h$ be $h-t$. Then $t\notin C$ and $2t\in C$.
Replacing those depths $2t$ by the absent halves $t$ leaves $q$ distinct
positive integers. If there are $v$ such defects,
\[
 e_0\ge q+\sum t\ge q+v\ge2v=2\delta_0-2.
\]
If there are $\ell\ge1$ low defects, their forced holes $2k$ have depths
at least $r-q$, spaced by at least two. Therefore
\[
 e_0\ge\ell(r-q)+\ell(\ell-1)+T_{q-\ell}-T_{q-1}.
\]
Subtracting $2q+2\ell$ from the right side gives
\[
 \ell(r-2q-2)+\tfrac32\ell(\ell-1)-q\ge0
\]
under $r\ge4q+1$. Since $\delta_0\le q+1+\ell$, again
$e_0\ge2\delta_0-2$. Finally, if $\mu>r-4$, then
$\Phi\le4q-(r-3)\le2$.
\end{proof}

\begin{lemma}\label{lem:threshold}
For a positive real set of size $r\ge14$, $\rho\ge0$.
If $r\ge6$ and its subset-sum support is not an arithmetic progression,
then $\rho\ge r-14$.
\end{lemma}
\begin{proof}
For $9\le i\le14$, put $q_i=\lfloor(i-1)/4\rfloor$. If $y_i\le q_i$,
\eqref{eq:small-inverse} and Lemma~\ref{lem:deleted} give $\Phi_i\le2$.
Otherwise $\Phi_i\le\Phi_{i-1}+1-q_i$. Starting at $\Phi_8\le10$ gives
successive bounds $9,8,7,6,4,2$. Larger increments $y_i\ge2$ preserve
the bound, and the audited cases $y_i=0,1$ do also.

For the second assertion, start with $\rho\ge-8$ at $r=6$ and induct.
If $y_r\le2$, \eqref{eq:small-inverse} gives at most two deletions, hence
$\delta_0\le4$, while stability gives $\mu\ge r-3$. Thus $\rho\ge r-9$.
If $y_r\ge3$ and the preceding support is not a progression, the defect
recurrence gives the result. If the preceding support is $t[0,L]_{\Z}$,
its next translate must be disjoint: any overlap would again produce a
progression. Consequently $y_r\ge(r-1)(r-2)/2$, and
$\rho_r\ge(r-1)(r-2)/2-10\ge r-14$ for $r\ge7$.
\end{proof}

\begin{lemma}[Span bound]\label{lem:span}
For a positive integer set $X$ of size $r\ge15$, with support hull-hole
count $h_X$, one has
\[
 \max X\le r+\rho(X)+h_X+2.
\]
\end{lemma}
\begin{proof}
Delete the largest element $r+q$ and let $\nu\ge0$ be the number of
previous defects it destroys. If primes denote the preceding prefix,
\[
 \rho(X)=\rho(X')+q-h_X+h_{X'}-2+2\nu.
\]
Now $\rho(X')\ge0$, $h_{X'}\ge0$, and $\nu\ge0$.
\end{proof}

\section{One common lattice and one exact budget}
Assume for now that $0\in A$ and $104\le n=|A|\le115$. Integralize
before using an integer inverse theorem. Form $B,D$ as in the covering
lemma and make two four-element peels $D\supset D_1\supset E$.
The elementary comparison of \eqref{eq:cover} and two copies of
\eqref{eq:peel} gives
\begin{equation}\label{eq:crude-core}
 s:=|\FS(E)|\le22d+2\Delta(D)+77\le24d+77.
\end{equation}
If $|B|\le7$, $D$ is one-sided and $E$ has at least $n-16\ge88$
nonzero elements of one sign. Then $s\ge1+T_{88}=3917$, contradicting
$s\le24\cdot115+77=2837$. Thus $|B|=8$ and
\[
 d=n-8,\qquad e=|E|=n-16,\qquad
 E=P\cup\{0\}\cup(-Q),\quad a=|P|,\ b=|Q|,\quad m=a+b=n-17.
\]

\subsection{Ensuring both signs and an opposite pair}
If $E$ contains $p$ pairs $\{x,-x\}$ and $q$ singly occurring absolute
values, then $e=2p+q+1$ and
\begin{equation}\label{eq:unpaired}
 |\FS(E)|\ge1+T_{p+q}+T_p=\frac{e^2+3+q^2}{4}.
\end{equation}
Indeed translate the support to replace negative elements by their positive
magnitudes. For ordered distinct magnitudes $b_i$, of multiplicity
$m_i\in\{1,2\}$, adjoining the $m_i$ copies of $b_i$ increases the
support by at least $i m_i$. The residues of $0,b_1,\ldots,b_{i-1}$
modulo $b_i$ are distinct, and each of these classes gains at least $m_i$
points. The smallest possible sum of these increments assigns multiplicity
two to the first $p$ positions, proving \eqref{eq:unpaired}.

Combining with \eqref{eq:crude-core}, where $d=e+8$, gives
\[
 q^2\le96e+1073-e^2,\qquad q\equiv e-1\pmod2.
\]
For $88\le e\le99$ this forces $p\ge23$ (see the table below).
Hence $a,b\ge23$, both peels really removed two elements of each sign,
and zero is not a diagonal exception in $D$.

\subsection{The exact slack identity}
Write $\mu=\mu(P)+\mu(Q)$ and $\rho=\rho(P)+\rho(Q)$. Set
\begin{align*}
 \Lambda&=1+T_a+T_b,&\kappa&=s-\Lambda-\mu,\\
 \theta&=\delta_0(P)+\delta_0(Q)+8-\Delta(D),&
 c&=2|2D|+18d+85-|\FS(D)|,\\
 \ell_1&=|\FS(D)|-|\FS(D_1)|-|R(D)|+1,\\
 \ell_2&=|\FS(D_1)|-s-|R(D_1)|+1,\\
 g&=4d-10-(|R(D)|-|R(D_1)|),&
 \T&=\kappa+\ell_1+\ell_2+g.
\end{align*}
All these slacks are nonnegative. For $\kappa$, use
$|U+V|\ge|U|+|V|-1$; for $g$, at most $4d-10$ pairs involve a
deleted point. The opposite pair just proved removes zero from the count
of possible diagonal exceptions. Direct cancellation gives the exact identity
\begin{equation}\label{eq:budget}
 \boxed{\rho+2\theta+c+\T=\Omega:=22d+97-\Lambda.}
\end{equation}

\subsection{Core normalization and propagation to the peeled layers}
We use the following sufficient form of Stanchescu's inverse theorem
\cite[Corollary 2]{Stanchescu}. If integer sets $U,V$, each of size at
least two, satisfy
\[
 |U+V|<|U|+|V|+\min(|U|,|V|)-3,
\]
their hull cardinalities on a common difference lattice are at most
$|U+V|-|V|+1$ and $|U+V|-|U|+1$, respectively.
For every possible oriented $(a,b)$ with $a,b\ge23$, $a+b=m$ and
$\Omega\ge0$, the table below verifies
\begin{equation}\label{eq:inverse-margin}
 T_{\min(a,b)}-2-\Omega>0.
\end{equation}
Thus $\kappa\le\Omega<\min(|\FS(P)|,|\FS(Q)|)-3$, which is a
stronger sufficient hypothesis. Apply the theorem to
$U=\FS(P)$ and $V=-\FS(Q)$ and divide by the common step. Because both
supports contain zero, $P,Q,E$ are now integral. If $h_P,h_Q$ and $z$
denote their support hull-hole counts, respectively, then
\begin{equation}\label{eq:holes}
 h_P,h_Q\le\kappa,\qquad z=h_P+h_Q-\kappa,\qquad
 0\le z\le\min(h_P,h_Q)\le\kappa.
\end{equation}
Moreover $h_X<|\FS(X)|-1$ for either side.

Let $R_p=\max P$ and $R_q=\max Q$. Lemma~\ref{lem:span} gives
$R_p+R_q\le m+\rho+\kappa+z+4$. Since
$|R(E)|\le2(R_p+R_q)+1$ and the next four points create at most $4e+6$
restricted sums, with $C_*=2m+4e+14-\Lambda$ we have
\begin{align}
 |\FS(D_1)|-2s
 &\le C_*+2\rho-\mu+\kappa+2z+\ell_2\notag\\
 &\le C_*+3\Omega.\label{eq:inner}
\end{align}
The second inequality follows from $z\le\kappa$ and
$\rho+\kappa+\ell_2\le\Omega$. The table verifies that it is strictly
negative. Any added point off the core lattice would give two disjoint
copies of $\FS(E)$ in $\FS(D_1)$, a contradiction. Hence $D_1$ is integral.
Similarly,
\begin{equation}\label{eq:outer}
 |\FS(D)|-2|\FS(D_1)|\le4d-10+\ell_1-\ell_2-s
 \le4d-10+\Omega-\Lambda<0,
\end{equation}
so $D$ is integral. No integrality of the outer eight points is assumed yet.

In the table, ``paired minimum'' is the lower bound on $p$ from
\eqref{eq:unpaired}, ``inverse'' is the minimum in
\eqref{eq:inverse-margin}, ``inner'' is the maximum of $C_*+3\Omega$,
and ``outer'' is the minimum of $\Lambda-(4d-10+\Omega)$.
All oriented pairs, including both reflections, are checked.

\begin{center}\small
\begin{tabular}{rrrrrr}
\toprule
$n$ & oriented pairs & paired minimum & inverse & inner & outer\\\midrule
115 & 1 & 36 & 1223 & -1845 & 2033\\
114 & 10 & 34 & 981 & -1721 & 1961\\
113 & 15 & 33 & 854 & -1597 & 1889\\
112 & 18 & 32 & 770 & -1477 & 1819\\
111 & 21 & 30 & 695 & -1357 & 1749\\
110 & 22 & 29 & 643 & -1241 & 1681\\
109 & 25 & 28 & 581 & -1125 & 1613\\
108 & 26 & 27 & 535 & -1013 & 1547\\
107 & 29 & 26 & 486 & -901 & 1481\\
106 & 30 & 25 & 446 & -793 & 1417\\
105 & 31 & 24 & 408 & -685 & 1353\\
104 & 34 & 23 & 371 & -581 & 1291\\
\bottomrule
\end{tabular}
\end{center}


The two extreme translates of $\FS(E)$ inside $\FS(D_1)$ overlap by
\eqref{eq:inner}; their hulls cover the new hull, leaving at most $2z$
holes. If their displacement is $u$, their union has at least $s+u-z$
points: its hull has $u$ more lattice points than the old hull and at most
$2z$ holes. Comparing this to the exact definition of $\ell_2$ gives
$\eta_2\le z+\ell_2$, where $\eta_2$ counts holes of $R(D_1)$ in
its integer hull. The same argument at the other peel, using
\eqref{eq:outer}, gives
\begin{equation}\label{eq:propagation}
 w:=\operatorname{holes}(\FS(D))\le4z,\qquad
 \eta_1\le2z+\ell_1,\qquad\eta_2\le z+\ell_2.
\end{equation}

\section{The four-extreme defect budget}
\begin{lemma}\label{lem:four}
Let $F=\{x_1<\cdots<x_v\}\subset\Z$, $v\ge2$, and add four integers
$a_1<a_2<x_1<x_v<b_2<b_1$ to obtain $D$. If $\eta_F,\eta_D$ count
restricted-sum hull holes and
$g_H=4v+6-(|R(D)|-|R(F)|)$, then
\[
 \eta_D+\eta_F+g_H\ge2v-4.
\]
\end{lemma}
\begin{proof}
Put $L=x_1+x_2$, $U=x_{v-1}+x_v$,
$h_+=b_1+b_2-U$, $h_-=L-a_1-a_2$, and $h=h_++h_-$.
Let $W_+$ be the union of the portions of $b_1+F,b_2+F$ above $U$.
Write $u=b_2-x_v$, $v_0=b_1-b_2$, $e_0=x_v-x_{v-1}$, and
$I_+=(U,b_1+x_v]_{\Z}$. If $H_+=|I_+|-|W_+|$, then
\[
 |W_+|\le u+1+v,\quad |W_+|=u+v_0+e_0-H_+,\quad
 2|W_+|\le h_++v+1-H_+.
\]
For the first bound, the terms of $b_2+F$ above $U$ use at most $u$
integers up to $x_{v-1}$ and the one point $x_v$; the other translate
has at most $v$ terms. Reflection gives the corresponding $W_-,H_-$.
Apart from $W_+,W_-$ and the two new extreme sums, at most four mixed
pairs $a_i+b_j$ create new values outside $[L,U]$. Let $q_0$ count
those not already counted. They lie in the unused parts of $I_+,I_-$,
so $q_0\le4$ and $q_0\le H_++H_-$. Therefore
\[
 2N_{\rm out}\le h+2v+6-(H_++H_-)+2q_0\le h+2v+10.
\]
If $N=|R(D)|-|R(F)|$, then $N\le N_{\rm out}+\eta_F$ and
$\eta_D=h+\eta_F-N$. Substitution yields
\[
 \eta_D+\eta_F+g_H=h+2\eta_F+4v+6-2N\ge2v-4.
\]
\end{proof}
Apply the lemma with $F=D_1$, $v=d-4$. Equations
\eqref{eq:holes} and \eqref{eq:propagation} imply, for $G=2d-12$,
\begin{equation}\label{eq:Tlower}
 \boxed{\T\ge\kappa+G-3z\ge G-2z,\qquad
 \T\ge z,\qquad\T\ge t_0:=\lceil G/3\rceil.}
\end{equation}
If $z>0$, both $h_P,h_Q$ are positive but smaller than their support
cardinalities minus one. Thus neither support is an arithmetic progression:
step one has no holes and step at least two has too many. Stability and
Lemma~\ref{lem:threshold} now give
\begin{equation}\label{eq:nonap}
 \mu(P)\ge a-3,\quad\mu(Q)\ge b-3,\quad\rho\ge m-28.
\end{equation}

\section{Twenty-four effective hole witnesses}
Only the hole threshold 14 is needed. We give its complete quantitative
proof to avoid any unspecified large-size assumption.

\begin{lemma}[Counted endpoint holes]\label{lem:counted}
If $i\ge2e_0+5$ and $y_i=e_0\ge0$, then, up to positive scaling,
\[
 X_i=[1,i+e_0]_{\Z}\setminus D_0,\quad |D_0|=e_0,
\]
and its support is
\[
 \mathcal S(L,H):=[0,L]_{\Z}\setminus(H\cup(L-H)),\quad
 H\subseteq[1,2e_0]_{\Z},\quad |H|\le e_0.
\]
\end{lemma}
\begin{proof}
The equality case handles $e_0=0$. Otherwise
\eqref{eq:small-inverse} first gives $\delta\le e_0$ deletions and maximum
$i+\delta$. Remove the maximum. Lev's theorem \cite[Theorem 1]{Lev}
states that for $B\subseteq[1,L_0]_{\Z}$, $|B|=s_0$, and
$L_0\le3s_0/2-2$,
\[
 [2L_0-2s_0+1,\ \sum B-(2L_0-2s_0+1)]_{\Z}\subseteq\FS(B).
\]
Apply this with $s_0=i-1$, $L_0=i+\delta-1$; the density condition is
$i\ge2\delta+5$. By symmetry its support has endpoint holes in
$[1,2\delta]$. Its total is at least $T_{i-1}>i+5\delta$. Adding
$i+\delta>2\delta$ overlaps the central intervals and keeps precisely
the same endpoint holes. Its support increment is $i+\delta$, so
$\delta=e_0$. Every lower hole is absent from the coefficient set itself;
there are only $e_0$ missing coefficients. Hence $|H|\le e_0$.
\end{proof}

\begin{lemma}[Effective hole test]\label{lem:effective}
For an integer $C\ge0$, define
\begin{gather*}
 K=\max\{k\ge3:k(k-3)/2\le C+1\},\qquad D_*=C+K,
 \qquad E_*=D_*+1,\\
 M=\max\left\{14,1+\left\lfloor
 \frac{1+\sqrt{9+8(E_*+42)}}2\right\rfloor\right\},
 \qquad J=\max\{K+1,19,M\}.
\end{gather*}
If
\begin{equation}\label{eq:hole-tests}
 7(r-J+1)>D_*,\qquad T_{J-1}-1>D_*,
\end{equation}
a positive $r$-element set whose support has, on a specified lattice,
$14<h<|\FS|-1$ must satisfy $\mu\ge r+C+1$.
\end{lemma}
\begin{proof}
Assume $\mu\le r+C$ and divide by the specified lattice step.
The set is integral and primitive: a gcd at least two would create at
least $|\FS|-1$ holes. Its support is not a progression.
If no prefix of size at least three has progression support, then
$y_j\ge1$ for $j\ge4$ and $\sum_{j=4}^r(y_j-1)\le C+3\le D_*$.
Otherwise let $k\ge3$ be the last progression prefix. The next support
translate is disjoint, giving $y_{k+1}\ge T_k-k$. Subsequent increments
are positive, so
\[
 \mu\ge r-1+k(k-3)/2,\qquad k\le K.
\]
In either case a tail starting no later than $K+1$ satisfies
\begin{equation}\label{eq:tail}
 y_j\ge1,\qquad \sum(y_j-1)\le D_*.
\end{equation}
The first test in \eqref{eq:hole-tests} forces some $i\in[J,r]$ with
$1\le y_i\le7$. Lemma~\ref{lem:counted} gives a scaled deleted interval
with endpoint state $H\subseteq[1,14]$, $|H|\le7$. Its coefficient set
has gcd one, so its scale $t$ is a positive integer. If $t>1$, primitivity
of the full set gives a first later coefficient off $t\Z$. At that step
the supports are disjoint and
\[
 y_j-1\ge T_{j-1}-j\ge T_{i-1}-1\ge T_{J-1}-1>D_*,
\]
contradicting \eqref{eq:tail}. Thus $t=1$.

Every later excess is at most $E_*$. We verify that the endpoint state
persists. For a current prefix of size $v\ge M$ with support
$\mathcal S(L,H)$, $L\ge T_v$, and its next element $x>v\ge14$.
If $x\ge L-28$, the portion of $x+\mathcal S(L,H)$ above $L$ has at
least $L-42\ge T_v-42>v+1+E_*$ points, contradicting its allowed
increment. The last strict inequality is exactly the definition of $M$.
Hence $14<x<L-28$. The two central integer intervals overlap or abut;
the lower holes remain $H$ and symmetry gives upper holes $L+x-H$.
Induction leaves at most $2|H|\le14$ holes in the final support, a contradiction.
\end{proof}

The following witnesses are sufficient, with no assertion of optimality.
The last columns are $7(r-J+1)-D_*$ and $T_{J-1}-1-D_*$, respectively.
The program also verifies $T_M-42>M+1+E_*$ using integer arithmetic.

\begin{center}\small
\begin{tabular}{rrrrrrr}
\toprule
$r$ & $C$ & $K$ & $D$ & $J$ & window margin & off-scale margin\\\midrule
32 & 83 & 14 & 97 & 19 & 1 & 73\\
33 & 89 & 15 & 104 & 19 & 1 & 66\\
34 & 96 & 15 & 111 & 19 & 1 & 59\\
35 & 102 & 15 & 117 & 19 & 2 & 53\\
36 & 109 & 16 & 125 & 19 & 1 & 45\\
37 & 110 & 16 & 126 & 19 & 7 & 44\\
38 & 116 & 16 & 132 & 20 & 1 & 57\\
39 & 122 & 17 & 139 & 20 & 1 & 50\\
40 & 128 & 17 & 145 & 20 & 2 & 44\\
41 & 129 & 17 & 146 & 21 & 1 & 63\\
42 & 135 & 18 & 153 & 21 & 1 & 56\\
43 & 142 & 18 & 160 & 21 & 1 & 49\\
44 & 147 & 18 & 165 & 21 & 3 & 44\\
45 & 149 & 18 & 167 & 22 & 1 & 63\\
46 & 155 & 19 & 174 & 22 & 1 & 56\\
47 & 162 & 19 & 181 & 22 & 1 & 49\\
48 & 167 & 19 & 186 & 22 & 3 & 44\\
49 & 168 & 19 & 187 & 23 & 2 & 65\\
50 & 175 & 20 & 195 & 23 & 1 & 57\\
51 & 182 & 20 & 202 & 23 & 1 & 50\\
52 & 187 & 20 & 207 & 23 & 3 & 45\\
53 & 188 & 21 & 209 & 24 & 1 & 66\\
54 & 195 & 21 & 216 & 24 & 1 & 59\\
55 & 202 & 21 & 223 & 24 & 1 & 52\\
\bottomrule
\end{tabular}
\end{center}


\section{A covering certificate retaining the core masses}
This section includes centers outside the target hull. Their errors are
absorbed by the available source span, so no additional case enumeration
for such centers is needed.
Write the eight outer elements as
\[
 p_1>p_2>p_3>p_4>M_p:=\max D,\qquad
 q_1>q_2>q_3>q_4>M_q:=-\min D.
\]
Let $S_p,S_q$ be the sums of the positive and negative magnitudes in $D$.
Put
\begin{gather*}
 \varepsilon_p=\sum P-T_a=\mu(P)+h_P,\qquad
 \varepsilon_q=\sum Q-T_b=\mu(Q)+h_Q,\\
 F=S_p+S_q+c-2(M_p+M_q)-2(R_p+R_q),\\
 F_A(a,b)=\tfrac12T_a+T_b+6a+8b+85,\\
 F_B(a,b)=\tfrac23T_a+T_b+6a+\tfrac{22}3b+\tfrac{236}3.
\end{gather*}
Every sum involving a core point, and every mixed-sign sum from $D\setminus E$,
lies in $[-M_q-R_q,M_p+R_p]$. Only ten positive-positive and ten
negative-negative pairs of the four peeled points per sign can lie outside.
Thus $|2D|\le M_p+M_q+R_p+R_q+21$. Using the definition of $c$ and
$|\FS(D)|=S_p+S_q+1-w$ gives
\begin{equation}\label{eq:Fupper}
 \boxed{F\le18d+126+w\le18d+126+4z.}
\end{equation}

\begin{lemma}[All ordinary units active]\label{lem:active}
If $c<d$, then after a possible reflection one of the following holds:
\begin{align}
 F&\ge F_A(a,b)+\tfrac12\varepsilon_p+\varepsilon_q,\label{eq:FA}\\
 F&\ge F_B(a,b)+\tfrac23\varepsilon_p+\varepsilon_q.\label{eq:FB}
\end{align}
There is no hypothesis limiting the number of out-of-hull centers.
\end{lemma}
\begin{proof}
Define, on the positive side,
\[
 \lambda_p=p_1-p_2-p_3-p_4,\quad
 \alpha_p=p_2-p_3-p_4,\quad\beta_p=2p_2-p_1-p_3-p_4.
\]
The ten ordinary centers are $\lambda_p+p_i$ ($1\le i\le4$),
$\lambda_p-q_j$ ($1\le j\le4$), $\alpha_p$, and $\beta_p$.
Reflection defines the negative side, represented as positive targets
and translates of $-D$. Pair $\lambda_p-q_1$ with $\lambda_q-p_1$.
If $\beta_p+\beta_q\le0$ (branch A), pair $\beta_p$ with $\beta_q$;
otherwise (branch B), pair $\lambda_p-q_2$ with $\lambda_q-p_2$.
These are the two capacity-$d$ pairings in \eqref{eq:cover}.

If any of the eighteen units were empty, it would cost $d$ in the cover
slack, contrary to $c<d$. Every unpaired center therefore contributes to
the actual support. The first pair's centers sum to less than
$-(M_p+M_q)$, so at most one even meets its open half-line. In branch A
the two unpaired centers $\lambda_p-q_2,\lambda_q-p_2$ both contribute,
and their sum is greater than $-(M_p+M_q)$. The identity
\[
 \beta_p+\beta_q=-(\lambda_p-q_2+\lambda_q-p_2)
                  -2(p_3+p_4+q_3+q_4)
\]
then makes $\beta_p+\beta_q<-(M_p+M_q)$. In branch B the second pair's
sum is already less than $-(M_p+M_q)$ by the same identity. Thus each
paired unit has precisely one contributing member.

All four own centers $b_i=\lambda_p+p_i$ contribute, and hence are
integers because $D$ and the target support are integral. Since
$\lambda_p=-b_1+b_2+b_3+b_4$, all $p_i$ are integral. Reflection proves
the same for all $q_i$.

Reflect so the first pair's contributor is $A_0=\lambda_p-q_1$.
In branch A put $B_0=\lambda_q-p_2$, which contributes. Then
\[
 \beta_p=-(A_0+B_0)-2(p_3+p_4)-q_2-q_3-q_4<-M_p,
\]
so the second contributor must be $\beta_q$. In branch B, $\beta_p$
is unpaired and contributes, giving
\[
 p_2-q_1=A_0+\beta_p+2p_3+2p_4>2M_p.
\]
It follows that $\lambda_q-p_2<-M_q$, and the second contributor is
$\lambda_p-q_2$. This proves that the following two combinations exhaust
the actual configurations.

For each selected slot, let $l_v,u_v$ count the points of its translate
below 1 and above its target maximum $S_p$ or $S_q$. Hull counts dominate
actual target counts; the cover, with 84 fixed contributions, gives
\begin{equation}\label{eq:loss}
 \sum_v(l_v+u_v)\le c.
\end{equation}
Define the nonnegative span allowances
\[
 U=M_p-a-4,\quad V=M_q-b-4,\quad q_p=R_p-a\le U,\quad q_q=R_q-b\le V.
\]
A contributing positive center obeys
\begin{equation}\label{eq:quantile}
 v\ge b+5-l_v-U\mathbf1_{v<0},\qquad
 v\le S_p-a-4+u_v+V\mathbf1_{v>S_p}.
\end{equation}
For $0\le v\le S_p$, count the $b+4$ negative elements and zero at the
lower end, and the $a+4$ positive elements at the upper end. If $v<0$,
contribution implies $-v<M_p$; among $[1,-v]$ at least $-v-U$ integers
belong to the positive part of $D$. This proves the corrected lower
bound. If $v>S_p$, contribution implies $v-S_p<M_q$; the corresponding
count of negative magnitudes proves the upper bound. The negative-side
versions interchange $a,b$ and $U,V$.

In branch A use the lower centers
\[
 \lambda_p-q_1,\quad\lambda_q-p_2,\quad\lambda_q-p_3,
 \quad\alpha_p,\quad\beta_q
\]
with coefficient one and the upper center $O=\lambda_p+p_1$ with
coefficient $1/2$. Their exact difference is
\[
 \tfrac12 O-\bigl((\lambda_p-q_1)+(\lambda_q-p_2)
 +(\lambda_q-p_3)+\alpha_p+\beta_q\bigr)
 =\tfrac12p_2+\tfrac52p_3+\tfrac32p_4+3q_3+3q_4.
\]
All losses used are distinct formal slots with coefficients at most one,
so \eqref{eq:loss}, \eqref{eq:quantile}, and integral outer spacings give
\begin{equation}\label{eq:Amid}
 \tfrac12S_p+c\ge\tfrac92M_p+6M_q+\tfrac72a+2b+44-\xi_A,
 \qquad\xi_A\le2U+\tfrac72V.
\end{equation}
The corrections count two positive lower rows, three negative lower rows,
and half a positive upper row.

In branch B use lower centers
$\lambda_p-q_1,\lambda_q-p_3,\lambda_q-p_4,\beta_q$ with coefficient
one, $\beta_p$ with coefficient $1/3$, and the upper center $O$ with
coefficient $1/3$. Here the difference is
$2p_3+2p_4+3q_3+3q_4$, and the result is
\begin{equation}\label{eq:Bmid}
 \tfrac13S_p+c\ge4M_p+6M_q+\tfrac{10}3a+\tfrac43b+38-\xi_B,
 \qquad\xi_B\le\tfrac43U+\tfrac{10}3V.
\end{equation}

Finally retain the exact core masses:
\[
 S_p\ge T_a+\varepsilon_p+3R_p+6+M_p,\qquad
 S_q\ge T_b+\varepsilon_q+3R_q+6+M_q.
\]
Substitute these into the unused $S_p/2+S_q$ in branch A and
$2S_p/3+S_q$ in branch B. After subtracting
$2(M_p+M_q)+2(R_p+R_q)$, the constant parts are exactly $F_A,F_B$.
The residual span terms, after all corrections, are at least
\begin{align*}
 3U+5V-\tfrac12q_p+q_q-\xi_A&\ge\tfrac12U+\tfrac32V\ge0,\\
 \tfrac83U+5V+q_q-\xi_B&\ge\tfrac43U+\tfrac53V\ge0.
\end{align*}
This proves \eqref{eq:FA} and \eqref{eq:FB}, even when several centers
are outside their target hulls. Formal-slot multiplicities are retained
throughout, so coincidences between numerical centers do not invalidate
the loss accounting.
\end{proof}

\section{Exact exclusion of all remaining scalar states}
There are finitely many necessary integer parameters, not finitely many
candidate real sets. For each $104\le n\le115$, enumerate all oriented
integers $a,b\ge23$, $a+b=n-17$, $\Omega\ge0$. There are 262 such
orientations in total. The preceding proofs apply to every actual set
with these parameters.

\subsection{At most fourteen core holes}
For each $z\in\{0,\ldots,14\}$ define
\[
 \rho_0=\begin{cases}0&z=0,\\m-28&z>0,\end{cases}\qquad
 c_{\max}=\Omega-\rho_0-\max\{t_0,z,G-2z\}.
\]
If $c_{\max}<0$, the state is impossible. Otherwise the program checks
$c_{\max}<d$, licensing Lemma~\ref{lem:active}. For $z=0$ use
$\varepsilon_p,\varepsilon_q\ge0$; for $z>0$ use
\[
 \varepsilon_p\ge a-3+z,\qquad\varepsilon_q\ge b-3+z.
\]
It then checks that both lower bounds in \eqref{eq:FA}, \eqref{eq:FB}
strictly exceed $18d+126+4z$. All 416 surviving scalar states fail in
this way. These are necessary relaxations; they need not be realizable.

\subsection{At least fifteen core holes}
If $z\ge15$, first require $\Omega-(m-28)\ge t_0$, or the budget is
already impossible. The 113 orientations surviving this condition have
both side sizes in $[32,55]$, exactly within the explicit witness table.
For either side put
\[
 u(r)=r+C(r)+1,\qquad v(r)=\max\{r-14,u(r)+2-2r\}.
\]
Because $h_X\ge z>14$ and $h_X<|\FS(X)|-1$, Lemma~\ref{lem:effective}
gives $\mu(X)\ge u(r)$, and $\delta_0(X)\le r$ gives $\rho(X)\ge v(r)$.
Let $H=\Omega-v(a)-v(b)$. Then
\[
 \T+c\le H,\qquad z\le\T\le H,\qquad c\le H-t_0.
\]
If $H<t_0$, discard the state. Only six orientations remain. In every
one, $H-t_0<d$. Use $\varepsilon_p\ge u(a)+z$ and
$\varepsilon_q\ge u(b)+z$ in the active certificate. The two margins over
\eqref{eq:Fupper} are bounded below by
\begin{align*}
 F_A(a,b)+\tfrac12u(a)+u(b)-(18d+126)-\tfrac52H,\\
 F_B(a,b)+\tfrac23u(a)+u(b)-(18d+126)-\tfrac73H.
\end{align*}
Both are strictly positive for every remaining orientation.

\subsection{Results and the real-set conclusion}
In the table below, a dash means that the preceding budget already
discarded all states. A reported margin is the minimum of both
certificate margins over all states in that column. All are exact
rational numbers. The $c$ columns give the largest remaining upper bound
on $c$, which is strictly below $d=n-8$.

\begin{center}\small
\begin{tabular}{rrrrrrr}
\toprule
$n$ & low states & $c_{\rm low}$ & low margin & high states & $c_{\rm high}$ & high margin\\\midrule
115 & 0 & $\text{--}$ & $\text{--}$ & 0 & $\text{--}$ & $\text{--}$\\
114 & 0 & $\text{--}$ & $\text{--}$ & 0 & $\text{--}$ & $\text{--}$\\
113 & 0 & $\text{--}$ & $\text{--}$ & 0 & $\text{--}$ & $\text{--}$\\
112 & 0 & $\text{--}$ & $\text{--}$ & 0 & $\text{--}$ & $\text{--}$\\
111 & 0 & $\text{--}$ & $\text{--}$ & 0 & $\text{--}$ & $\text{--}$\\
110 & 0 & $\text{--}$ & $\text{--}$ & 0 & $\text{--}$ & $\text{--}$\\
109 & 0 & $\text{--}$ & $\text{--}$ & 0 & $\text{--}$ & $\text{--}$\\
108 & 0 & $\text{--}$ & $\text{--}$ & 0 & $\text{--}$ & $\text{--}$\\
107 & 9 & $18$ & $\frac{545}{2}$ & 0 & $\text{--}$ & $\text{--}$\\
106 & 42 & $43$ & $209$ & 0 & $\text{--}$ & $\text{--}$\\
105 & 151 & $68$ & $168$ & 0 & $\text{--}$ & $\text{--}$\\
104 & 214 & $92$ & $\frac{263}{2}$ & 6 & $7$ & $351$\\
\bottomrule
\end{tabular}
\end{center}


This excludes all zero-containing feasible sets of sizes 104 through
115. Any feasible real set has size at most 115 by the audited foundation.
If one of size at least 104 omitted zero, adjoining zero would either
contradict that foundation or give one of the excluded zero-containing
sizes. This proves Theorem~\ref{thm:main}. No step assumes that feasibility
passes to arbitrary subsets.

\section{Reproduction, audit boundary, and the next obstruction}
The new verifier is \texttt{verification/upper103.py}, using Python 3.10+
and only the standard library. From the repository root run:
\begin{verbatim}
python verification/upper103.py \
  --output results/2026-10-14/upper103/result.json
\end{verbatim}
It reconstructs the four center identities over the rationals, checks
span absorption, all 24 effective-hole witnesses and all strict gates,
all lattice margins, all low-hole states and all high-hole orientations.
Integer square roots and rational arithmetic are used throughout.
No floating-point tolerance, external solver, archived endpoint certificate,
or claimed historical PASS result is an input. Checks use explicit
exceptions and remain enabled under \texttt{python -O}.

The output is an arithmetic audit of the displayed reductions. It is
not a formal proof-assistant certification of the analytic lemmas or of
the cited theorems. Its two deliberately invalid effective-hole rows are
rejected. It also records the obstruction at the next endpoint:
\[
 n=103,\quad a=b=43,\quad z=0,\quad d=95,\quad
 \Omega=294,\quad G=178,\quad c\le116.
\]
Here $116\ge95$, so the proof does not force all ordinary units active.
This is a failure of the current relaxation, not a feasible 103-element
example. Continuing below 103 requires a stronger slack bound or a proof
covering at least one empty unit. The verifier deliberately does not
certify this next endpoint with the all-active lemma.

The threshold-14, non-progression, effective-hole, and active-center ideas
come from the archived continuation. This report rechecks and consolidates
the needed portion into one uniform descent, retains the exact core mass,
and makes the outside-center span absorption explicit. It leaves the
longer 103-to-84 chain outside the current independent audit.

\begin{thebibliography}{9}
\bibitem{CDK} R. Carpenter, C. Defant and N. Kravitz,
\emph{Sets with few subset sums}, arXiv:2605.05498v2 (2026),
Theorem 1.1 and Lemma 2.1.
\url{https://arxiv.org/html/2605.05498v2}.
\bibitem{BGGH} E. Balandraud, B. Girard, S. Griffiths and Y. Ould Hamidoune,
\emph{Subset sums in abelian groups}, European Journal of Combinatorics
\textbf{34} (2013), 1269--1286, Lemma 3.6 and Theorem 1.5.
\url{https://arxiv.org/abs/1112.1929}.
\bibitem{Stanchescu} Y. Stanchescu,
\emph{On addition of two distinct sets of integers},
Acta Arithmetica \textbf{75} (1996), 191--194, Theorem 4 and Corollary 2.
\url{https://doi.org/10.4064/aa-75-2-191-194}.
\bibitem{Lev} V. F. Lev, \emph{On consecutive subset sums},
Discrete Mathematics \textbf{187} (1998), 151--160, Theorem 1.
\url{https://doi.org/10.1016/S0012-365X(98)80006-X}.
\end{thebibliography}
\end{document}
